\subsection*{Síntese comparativa dos fluxos examinados}
\label{sec:comparacao-servicos}
A Tabela~\ref{tab:comparacao-servicos} reúne as evidências documentais descritas nesta seção. ``Observado'' significa visível nas capturas examinadas; ``não verificado'' indica insuficiência do acervo. Uma opção ausente do fluxo examinado não é declarada inexistente no serviço. Não foi realizado um ensaio simultâneo de usabilidade entre produtos.
\begingroup\small\SingleSpacing\setlength{\parskip}{0pt}\setlength{\tabcolsep}{4pt}
\begin{longtable}{>{\raggedright\arraybackslash}p{2.4cm}p{5.5cm}p{6.8cm}}
\caption{Serviços e continuidade da consulta: alcance da evidência.}\label{tab:comparacao-servicos}\\\hline
Serviço & Funcionalidade documentada & Limite da comparação\\\hline\endfirsthead
\hline Serviço & Funcionalidade documentada & Limite da comparação\\\hline\endhead
\hline\endfoot
QRCodeAdesivo & Página web contextualizada por parada, linhas e links individuais; Figura~\ref{fig:sptrans-parada}. & Seleção por destino ausente da página examinada. Leitura física pelo grupo e funcionamento completo não verificados; não se conclui encerramento do serviço.\\[5pt]
Olho Vivo & Lista de previsões por ponto e mapa da linha selecionada; Figuras~\ref{fig:olho-vivo-previsoes} e~\ref{fig:olho-vivo-mapa}. & Comparação espacial conjunta e planejamento por destino não observados nesse fluxo. É serviço distinto do QRCodeAdesivo.\\[5pt]
Cittamobi & Busca por destino, alternativas com conexões e mapas nas telas analisadas; Figura~\ref{fig:cittamobi-destino-rotas}. & Integração completa QR físico--destino no navegador não verificada. Destino, favoritos e conexões não são novidades do PaP.\\[5pt]
Santos / Quanto Tempo Falta & Resposta institucional reproduzida pelo grupo e descrição de vídeo relatam QR que abre página web da parada. & Evidência indireta: vídeo original não analisado; destino ausente apenas do percurso descrito. Data e setor da resposta ao protocolo 7657-2026 não informados.\\[5pt]
Google Maps e Moovit & Planejamento por destino e uso de favoritos registrados no estudo de ações em celular, em 25 e 26/10; Apêndice~\ref{ap:cliques}. & Destinos, anúncios e estados iniciais diferem. Entrada por QR físico não examinada; contagens não estabelecem superioridade geral.\\[5pt]
PaP & Entrada contextual por URL da parada, escolha do destino, alternativas, transferências e mapa no navegador; capturas e testes internos deste relatório. & Cadeia de navegação implementada; leitura física e benefício aos passageiros ainda não demonstrados. Não se estabelece superioridade entre serviços.\\
\end{longtable}
\fonte{Síntese do acervo examinado pelo grupo em 2026; capturas e procedência descritas nesta seção e no Apêndice~\ref{ap:cliques}.}
\endgroup
A contribuição técnica é a continuidade desse conjunto de ações em uma implementação web contextualizada. A hipótese de redução de esforço exige tarefas equivalentes, resultado correto e compreensão do percurso, não apenas a contagem de recursos disponíveis.
